\documentclass[leqno, a4paper, 12pt]{jsarticle}
\usepackage{amssymb}
\usepackage{amsthm}
\usepackage{amsmath}
\usepackage{comment}
\usepackage{url}

\usepackage{enumitem}
%\usepackage{showkeys}


%定理環境 thmはあまり使わずpropをなるべく使う。
%主張は基本的に証明の中で使い、そのため番号を付けない。
%注意環境を付けてみた。
\newtheorem{thm}{定理}
\newtheorem{prop}[thm]{命題}
\newtheorem{cor}[thm]{系}
\newtheorem{lem}[thm]{補題}
\newtheorem{df}[thm]{定義}
\newtheorem{exam}[thm]{例}
\newtheorem{fact}[thm]{事実}
\newtheorem*{claim}{主張}
\newtheorem*{cau}{注意}
\newtheorem*{rmk}{注意}
\renewcommand{\proofname}{証明}
%矢印たち。
%\toは\rightarrowと同じ。\getsは\leftarrowと同じ。同値は\iff
\newcommand{\To}{\Longrightarrow}
\newcommand{\Gets}{\LongLeftarrow}
%黒板太字
\newcommand{\nn}{\mathbb{N}}
\newcommand{\zz}{\mathbb{Z}}
\newcommand{\qq}{\mathbb{Q}}
\newcommand{\rr}{\mathbb{R}}
\newcommand{\cc}{\mathbb{C}}
%無限大
\newcommand{\mgn}{\infty}
%差集合は\setminusで統一する。等号の否定は\neq
%真部分集合は\subsetneqq　偏微分は\partial
%ディスプレイスタイル。\dfracでディスプレイ分数
\newcommand{\dpst}{\displaystyle}

\title{論文メモ
\large{\\--- ペアノ微分 ---}
}
\author{ナンブ　キトラ}
\date{\today}
\begin{document}
\maketitle

本棚を整理したいので，
溜まっている論文のメモを書いて未来への手紙とすることにする．

以下の論文の束は
ペアノ微分とかタイラー展開を勉強しようとして
そのままになっていたものだと思う．

$\rr$の適切な区間で定義された関数
$f$について$f^{(n)}$が
点$a$で存在するならば
\[
f(x)=\sum_{i=0}^{n}\frac{f^{(i)}(a)}{i!}h^{i}+o(h^{n})
\]
となるというのがペアノの剰余項というものである．
逆にこれを使って微分を定義するというのがペアノ微分である．
つまり点$a$における$f$のペアノ微分とは
$f_{1}(a)\dots f_{n}(a)$であって，
\[
f(x)=\sum_{i=0}^{n}\frac{f_{i}(a)}{i!}h^{i}+o(h^{n})
\]
を満たすもののことである．
$f^{n}$が存在するならば
$f_{n}=f^{(n)}$であるが逆は一般には成り立たずそこんところが魅力らしい．
またこの概念はWhitneyの微分可能性と
関係ありそうな気がするが，あの概念は各導関数が連続になるのでこの概念ともまた異なりそうだ．
\section{ペアノ微分}

1953年に
\cite{MR0062207}
では，通常の導関数$g$（$f^{\prime}=g$となる$f$が存在するということ， このような$g$は一般には連続とは限らない）
が以下のことを満たすと注意する．
\begin{enumerate}[label=\textup{(\arabic*)}]
\item 
$g$は第一級のBaire関数

\item 
$g$は中間値の定理の帰結を満たす．
\item 
$g$は平均値の定理を満たす(原始関数の$f$も含めた意味で)．

\item 
$A$, $B$は
$B<A$を満たすとする．
このとき
$\{\, x\mid B<g(x)<A\, \}$
は空であるかもしくは正の測度をもつ．
\end{enumerate}
これを踏まえた上で，
$f$の$n$次のペアノ微分$f_{n}$
も上と同様の性質を満たすことを
証明している．


1969年の論文
\cite{MR0233944}
ではペアノ微分と通常の導関数
の存在に関して何やら研究している．
Approximate derivativeも含めて研究している．
この概念はドメインをほとんど至る所埋め尽くす集合$E$
上で微分係数を計算した値で定義される．普通の微分は$E$じゃなくてドメイン全体で点を近づける．

1970年の論文
\cite{MR0259041}で
は
ちょっと歪みをかけたような$n$次の差分の極限が存在することと$n$次のPeano微分が存在することが同値であることを証明している．

1973 年の論文
\cite{MR0311845}
はちょっとよくわかんなかった．ごめん．
測度論的性質を研究している気がする．

1993 年の論文
\cite{MR1234991}
ではさらに一般化しててよくわかんない．
ごめん！

1997年の論文
\cite{MR1396976}
は
\cite{MR1234991}の発展だと思う．
path derivative というのを
研究している．

1998年の論文
\cite{MR1605380}
もよくわかんなかった．\cite{MR1234991}の発展だと思う．

2008 年の論文
\cite{MR2373608}
は多変数のペアノ微分に関して
それが実際に意味で微分可能であるかどうかの十分条件を与えているような気がする．
2013年の論文
\cite{MR3043022}
も多変数のPeano微分の話だと思う．

2013/14年の論文
\cite{MR3261899}
では完全集合$P$上で定義された関数の
ペアノ微分を考察している．
この研究の感じはWhitney微分を彷彿とさせるのでこれ面白んじゃないかな．
論文
\cite{MR3658803}もなんかそういう話で面白そう．
論文\cite{MR3261899}
では$P$上で定義された無限開微分可能関数で
$f(P)=P^{2}$
となるものを構成し，
連続微分可能な関数$f\colon \rr\to \rr^{2}$は
絶対に完全集合をその平方に写さないことを証明している．
なんかすごい面白い気がする．

\section{テイラー展開}


1937 
\cite{MR1523823}
テイラーの定理が証明されている．
いい感じの定数おいて最後に求めるやつ．
帰納的ににやるやつ．

1949
年
\cite{MR0028907}
なんか$f^{(k)}(a)$だけではなく
$f^{(k)}(x)$も使うテイラー展開を証明している．

1963
\cite{MR1532320}
クラスルームノート．
いろいろな形式のテイラー展開を解説している．
リマーク２では$[0, 1]$上の関数$f$で
$f^{\prime}$も常に存在するし，$f^{\prime\prime}(0)$
も存在するが
$f^{\prime}$がリーマン積分不可能な例を挙げている．



1967
\cite{MR1534508}
では
$f^{(n)}$が存在するための
必要十分条件をテイラー展開で特徴付けしている．

1980
\cite{MR0567920}
では
テイラーの定理に出てくる
$a+\theta(x-a)$の
$\theta$の$x\to a$の極限について調べている．
この極限についてはなんと先行研究があるらしい．
それを多変数に拡張している．
こういう極限を計算すると
数値計算へと応用があるようだ．

1989
\cite{MR1541471}
では多変数のテイラーの定理の短い証明を
与えている気がする．

1990
\cite{MR1096374}
では
多変数のテイラー展開をなるべく小さい
仮定で証明しているような気がする．

\section{その他}
Saksの論文
\cite{MR1574835}
はなんかもうよくわかんないんだけど，
3次の差分の極限のが正みたいな
仮定をして導関数の連続性を証明している．


1947
\cite{MR0019712}では以下の定理を証明している．
\begin{thm}
$f$は$[a, b]$上で定義された関数でこの上で微分可能とする．
このとき任意の$\alpha, \beta$に関して
集合
\[
\{\, x\in [a, b]\mid \alpha<f^{\prime}(x)<\beta\, \}
\]
は空集合であるか，もしくは正の測度を持つ．
ここで$f^{\prime}$は無限大の値も許すらしい．
\end{thm}
この性質を
Denjoy--Clarkson property 
と呼ぶらしい．

1962 
\cite{MR0151558}
では上の定理を
proximate derivative
にまで拡張している．


1970
\cite{MR0259040}
では上の定理を二次元に拡張してる．

%READ!
%myplain is the same to amsplain style. 
\nocite{*}
\bibliographystyle{myplaindoidoi}
\bibliography{pd.bib}



\end{document}